\Needspace{18\baselineskip}
\section{Worked problems}

\Needspace{12\baselineskip}
\subsection{Naive check budget}
\textbf{Problem.} \emph{Derivation.} With $V=200{,}000$, $B=64$ and an 8-ms model step, what
average $t_c$ would keep naive equation~\eqref{eq:structured-naive} below 10\% of the
step? Explain why the answer motivates precomputation.

\textbf{Solution.} Ten percent of 8 ms is 0.8 ms. Equation~\eqref{eq:structured-naive} gives
$t_c\le0.0008/(64\cdot200000)=6.25\times10^{-11}$ s, or 0.0625 ns/check.
That is not a realistic scalar parser budget; the vocabulary work must be avoided,
amortized, parallelized or overlapped.

\Needspace{12\baselineskip}
\subsection{Sparse versus dense}
\textbf{Problem.} \emph{Representation.} Compare bytes written by a dense bitset with a sparse
list when 12 of 200,000 token IDs are legal. State the ID width you assume.

\textbf{Solution.} A 200,000-bit dense mask is 25,000 bytes. Twelve legal 32-bit token IDs need only
48 bytes before container metadata, so sparse allowed IDs win strongly at that state.

\Needspace{12\baselineskip}
\subsection{Tool-call security}
\textbf{Problem.} \emph{Security.} Give an example of a syntactically valid tool call that must
still be rejected by authorization policy.

\textbf{Solution.} A grammar may accept a transfer tool with a numeric amount while authorization rejects
it because the caller lacks permission, the amount exceeds a limit, or confirmation is
missing. Syntax and authority are independent guarantees.

\Needspace{12\baselineskip}
\subsection{EOS}
\textbf{Problem.} \emph{Trace.} For the lab grammar, trace parser states after tokens
\texttt{\{"ok":}, \texttt{tr}, \texttt{ue\}} and identify where EOS first becomes
legal.

\textbf{Solution.} EOS is legal only after the recognizer reaches an accepting state representing the
complete closing brace; accepting a prefix such as \texttt{\{"ok":tr} would violate
the language even though it can still be extended to a valid string.
